\section{Introduction}
\label{sec:intro}

Rolling-element bearings fail in service, and the vibration signature of a spall or a crack is well understood: an impact
train at a kinematic frequency fixed by the bearing geometry and the shaft speed, modulated by load and smeared by cage slip
\cite{randall2011}. Condition-monitoring practice turns that signature into a decision --- replace the bearing, or leave it
--- and the value of an automated diagnosis is the reliability of that decision on a bearing the diagnostic system has not
seen before.

Data-driven diagnosis has moved towards transfer learning to meet that requirement. Unsupervised domain adaptation
transfers a model trained under one operating condition to another without target labels \cite{zhao2021udtl}, and source-free
domain adaptation (SFDA) goes further: during adaptation only the trained source model and unlabelled target recordings are
available, which suits an operator who cannot ship raw data between sites \cite{liang2020shot,sdalr2025}. Reported
accuracies are high. Vieira et al.\ sampled 18 supervised bearing-diagnosis papers published in 2025: every one reported
above \SI{94}{\percent}, five reported \SI{100}{\percent}, and none split training and test data by physical bearing ---
eight split at random, nine by condition (load, speed or noise level) and one did not state its split
\cite[Table~1]{vieira2026}. Source-free results sit at the same level: the method we re-evaluate, SDALR, reports
\SI{96.8}{\percent} on Paderborn \cite{sdalr2025} on a task whose eight classes are eight individual bearings, the same
specimens on the source and the target side.

A vibration dataset carries two kinds of information at once. One is the fault physics that should transfer: the impact
train at the outer-race or inner-race passing frequency. The other is everything particular to one specimen --- the extent
and shape of its damage, its mechanism, its mounting, clearance and resonances. When the same physical bearing appears in
training and in test, the particulars are sufficient to score well, and no experiment in that protocol can tell the two
apart. Splitting by physical bearing removes the shortcut, and in supervised settings it has been shown to change
conclusions \cite{hendriks2022,wheat2024leakage,vieira2026,kaya2026}; the same lesson was learned in clinical machine
learning, where record-wise rather than subject-wise splits overestimate accuracy \cite{saeb2017}.

SFDA deserves separate scrutiny for a reason specific to how it works. Adaptation is driven by the source model's own
predictions on target data --- pseudo-labels, centroids or neighbourhood structure --- refined over iterations. If the
source model has learned the particulars of its training specimens, adaptation might sharpen that error on a new bearing,
or it might repair it. Which one happens, how much performance is at stake, and what physics-based pseudo-label filters can
recover are empirical questions that a leaky protocol cannot ask. Physics-based filters are increasingly proposed:
characteristic-frequency evidence gates and weights pseudo-labels in the EAGLE module of Bio-SFDA \cite{yoon2026biosfda}
and builds physical pseudo-labels in PIUDA \cite{jia2024piuda}, and a physics consistency validator supervises confidence
in PCTL \cite{pctl2026}. These modules are judged by end accuracy; how often they declare a fault on a healthy bearing ---
which decides whether they are safe inside an adaptation loop --- is not reported in those works.

This paper follows one thread through those questions, with every decision rule fixed in a versioned protocol before the
corresponding run.

\begin{enumerate}
  \item \textbf{How much survives, and where the loss arises.} One source model is adapted to the same bearings and to
        unseen bearings at the same target condition (Section~\ref{sec:protocol}). We measure the gap on ten
        pre-registered random bearing splits, two mechanical folds and a second rig, for two SFDA methods and for
        non-adaptive baselines, before and after adaptation, and for each bearing when it is seen and when it is unseen
        (Sections~\ref{sec:res-collapse} and~\ref{sec:res-second}).
  \item \textbf{Why unseen bearings fail.} Probes, a per-bearing table drawn from the dataset's own damage profiles and a
        pre-registered source-diversity curve separate a fixed identity shortcut from under-sampled damage signatures;
        a classifier on kinematic envelope features tests whether the loss is specific to the input representation
        (Section~\ref{sec:res-mechanism}).
  \item \textbf{What any pseudo-label filter could recover.} An oracle filter that keeps only correct pseudo-labels gives a
        reference for keep-or-withhold filters (Section~\ref{sec:res-decomp}).
  \item \textbf{What physics gating actually buys.} Four gates are compared by false acceptance on 480 real healthy
        recordings and over their full operating curves, then inserted into the adaptation loop as paired arms
        (Section~\ref{sec:res-gating}).
\end{enumerate}

The novelty is narrow, and we state it against the nearest work. That bearing overlap inflates diagnosis accuracy is
established: Hendriks et al.\ on CWRU \cite{hendriks2022}, Wheat et al.\ across run, day and part splits
\cite{wheat2024leakage}, and Vieira et al., who also hold the training set fixed and change only the test set, including a
same-bearing, same-condition repetition-wise split on Paderborn \cite[Section~6.2]{vieira2026}. Knap et al.\ built a
recording-separated cross-domain benchmark with one cross-bearing scenario \cite{knap2026leakagesafe}, and concurrent
work adapts across Paderborn conditions with held-out bearings, without a leaky baseline or a source-free method
\cite{nagaswetha2026}. All of these are supervised, non-adaptive or source-dependent. What this paper adds is (i) the
fixed-source design applied to SFDA, with a source-only arm that locates the loss before adaptation and a within-bearing
contrast that holds the specimen fixed; (ii) a pre-registered source-diversity curve and a kinematic-feature control that
bear on why unseen bearings fail; (iii) the perfect-pseudo-label reference of Schlachter et al.\ \cite{schlachter2025},
applied as a keep-or-withhold oracle and set against the loss; and (iv) false-acceptance curves on real healthy recordings
for the physics gates that such methods insert, with the gates' coverage traced to the size of each specimen's damage.

We claim no new network, detector or adaptation algorithm. Statistical envelope thresholds are established
\cite{borghesani2013,randall2011} and harmonic-coordinate features close to the ones used here were published by Matania et
al.\ \cite{matania2025}. The contribution is knowledge about where the error of source-free bearing diagnosis lies, what
physics-based filtering can and cannot repair, and the evidence a result in this area should carry before it is believed.
